\section{Schinzel's Hypothesis H: the fixed-class implication}
\label{sec:schinzel}

Theorem~\ref{thm:classexist} supplies the canonical aligned-class
construction for every cell.  Section~\ref{sec:irreducibility} proves, for
every fixed nonzero rational $Z$, that its octic is irreducible in
$\Q(a)[b]$, and quantitative Hilbert irreducibility selects an irreducible
specialization inside the existing character-admissible $a$ progression.
The L20 admissibility package is therefore unconditional.  For any such
fixed class, classical Schinzel's Hypothesis~H \cite{schinzel} is the sole
remaining input and supplies the member.
% src: data/l22_elimination.jsonl (fixed-Z-theorem row)
% src: data/l20_admissible.jsonl; data/l18_schinzel_implies_h.jsonl

\subsection{The operational contract}

A place is recorded as an obstruction only when its valuation is odd
\emph{and} its Hilbert symbol is $-1$: the kernel appends a place under
exactly that conjunction, and the ladder returns \emph{zero} when every
remaining place is proved with symbol $+1$.
% src: l13_filter.py smooth_emergent / cofactor_decide
Emergent-freeness therefore means \emph{no place of symbol $-1$}, not
\emph{no odd valuation outside the frozen set}. This distinction is
load-bearing: $197$ of the $293$ grid closures have exactly one
odd-valuation outside prime whose symbol is forced $+1$ by reciprocity.
% src: THEOREMS.md L14/L17 (197 prime-rung closures, 197/197 singleton +1)
The stronger reading would not follow from prime values, since a prime
value deliberately creates one odd place.

\subsection{The theorem}

Fix a verified aligned class with data $(a,\varepsilon,f,q_1,N)$, put
$Q(t) = q_1 + Nt$, $b(t) = \varepsilon f Q(t)$, and let
$S = \{2,3,5,7\} \cup \supp(\alpha) \cup \supp(\delta) \cup \supp(f)$.
Write the composed class polynomial as
\[
   F(t) \;=\; P\bigl(\varepsilon f (q_1 + Nt)\bigr) \;=\; c\,G(t),
\]
with $c > 0$ and $G \in \Z[t]$ primitive with positive leading
coefficient. The relevant hypothesis on $c$ is that its \emph{square
class} is supported on $S$. Literal $S$-support is false in general: for
$(w,z) = (3, 3/11)$, $a = 1$, $q_1 = 19$ one gets
$c = 9072/15692141883605$, whose outside part is $11^{-12}$. Only the
square class matters, because an outside place of even valuation
contributes Hilbert symbol $+1$.
% src: data/l20_admissible.jsonl (normalization block, audited counterexample)

\begin{theorem}[Classical Schinzel H supplies emergent-free members; PROVED implication]
\label{thm:schinzel}
Assume the class certificate (all symbols at $S$, at the moving prime
$Q$, and at infinity equal $+1$ throughout the progression), assume the
square class of $c$ is supported on $S$ and $G(t)$ is a $p$-adic unit for
every $p\in S$, and assume $Q$ and $G$ are irreducible with positive
leading coefficients whose product has no fixed prime divisor.  Then
Schinzel's Hypothesis~H for $\{Q,G\}$ implies that the class contains
infinitely many members with ladder verdict \emph{zero}.
\end{theorem}

\begin{proof}[Proof sketch]
Schinzel supplies infinitely many $t$ with $Q(t)$ and $G(t)$
simultaneously prime; the side conditions (Cauchy bound, $b\ne0,1$, and
finitely many forbidden coincidences) exclude only finitely many $t$.
Put $R=G(t)$.  The part of $c$ outside $S$ need not be trivial, but every
one of its valuations is even.  At the moving prime, $b\equiv0\pmod Q$
and the kernel identity gives
$P(b)\equiv-\delta a^4Z^4A^2\pmod Q$, a unit, so $R\ne Q$.
Consequently $R$ is the unique additional outside place of odd valuation;
every other outside place has even valuation and $v_p(d)=0$, hence symbol
$+1$.  The certificate pins the symbols at $S$, at $Q$, and at infinity
to $+1$.  Hilbert reciprocity forces $(x_0,d_0)_R=+1$.  Thus no place
carries symbol $-1$, and the ladder verdict is \emph{zero}.
% src: data/l18_schinzel_implies_h.jsonl (fixed-class implication)
% src: data/l20_admissible.jsonl (content formula and square-class correction)
% src: THEOREMS.md L9 (reciprocity steering), h10q.py odd-prime symbol formula
\end{proof}

\subsection{Uniform admissibility status}

For the class of Theorem~\ref{thm:classexist}, let
$L(t)=q_1+Nt$ and normalize
$P(fL(t))=cG(t)$ with $c>0$ and $G\in\Z[t]$ primitive.  The exact status
is as follows.

\begin{table}[ht]
\centering
\small
\begin{tabular}{@{}p{.09\textwidth}p{.49\textwidth}p{.30\textwidth}@{}}
\toprule
condition & requirement & uniform status\\
\midrule
(a) & irreducible polynomials with positive leading coefficients &
\textsc{proved} for every fixed $Z\in\Q^\times$; Hilbert selection in the
admissible $a$ progression \textsc{proved}\\
(b) & $\gcd(q_1,N)=1$ & \textsc{proved}\\
(c) & $L(t)G(t)$ has no fixed prime divisor & \textsc{proved}\\
(d) & $G(t)$ is a $p$-adic unit at $S$ and the frozen square classes are
constant & \textsc{proved}\\
\bottomrule
\end{tabular}
\caption{Admissibility of the uniformly constructed square-branch classes.}
\label{tab:admissibility}
\end{table}
% src: data/l20_admissible.jsonl; data/l22_elimination.jsonl (fixed-Z-theorem row)

For (a), the constant and leading coefficients of $P$ are
$P_0=P_8=4a^8AZ^4>0$, and the affine substitution
$b=f(q_1+Nt)$ preserves factorization over $\Q$.  Thus $G$ has positive
lead and
\[
 G\text{ irreducible over }\Q
 \quad\Longleftrightarrow\quad
 P\text{ irreducible over }\Q.
\]
Theorem~\ref{thm:fiber-irred} proves the vertical polynomial irreducible
for every fixed $Z\in\Q^\times$, and
Corollary~\ref{cor:admissible-irred} selects a good $a$ in the required
character progression.  Thus the fixed cell, not merely the generic
two-variable family, is licensed.
% src: data/l20_admissible.jsonl; data/l22_elimination.jsonl

For (b), the prime support of $N$ is exactly $S$, while the CRT base
$q_1$ is a unit modulo every prime in $S$.  For (c), primes in $S$ never
divide either value by the Taylor congruences.  If $p\notin S$, the linear
factor has one root modulo $p$; a fixed divisor would force the degree-$8$
polynomial $G$ to vanish at the other $p-1$ residues, hence $p\le9$.
Every such prime already lies in $S$, a contradiction.  For (d), the
Taylor-exponent construction gives coefficientwise
\[
 G(t)\equiv G(0)\not\equiv0\pmod p
 \quad(p\in S\text{ odd}),\qquad
 G(t)\equiv G(0)\not\equiv0\pmod8
\]
and the analogous congruences for $L(t)$.  The resulting ratios are local
squares, so every frozen Hilbert symbol is constant for all integer $t$.
% src: data/l20_admissible.jsonl (uniform proofs for conditions (a)--(d))

The modular engine remains an exact finite cross-check.  The L20 class
audit certified $353/353$ independently reconstructed grid classes and
$706/706$ legacy cell--branch rows, while the L21 audit certified all
$353/353$ recorded target fibers.  The L22 endpoint--Newton theorem is
strictly stronger in scope: it proves the vertical statement for every
fixed nonzero rational $Z$, without extrapolating from those finite rows.
% src: data/l20_admissible.jsonl; data/l21_irred_generic.jsonl
% src: data/l22_elimination.jsonl (fixed-Z-theorem row)

A separate $24$-row cross-check re-ran the proven-primality kernel end to
end: every sampled member had one proved-prime ladder residual with symbol
$+1$ and verdict \emph{zero}.
% src: data/l18_schinzel_implies_h.jsonl (24 verified_prime_rung rows)

\begin{remark}[Exact scope]
Schinzel supplies members only after a chosen class satisfies every row of
Table~\ref{tab:admissibility}.  Every one of those rows is now PROVED on the
fixed canonical branch: Theorem~\ref{thm:fiber-irred} and
Corollary~\ref{cor:admissible-irred} supply irreducibility, while L20
supplies normalization, positivity, coprimality, absence of a fixed
divisor, and frozen local data.  Therefore the implication has the exact
form
\[
 \text{classical Schinzel H}
 \quad\Longrightarrow\quad
 \text{the conditional six-quantifier conclusion}.
\]
Schinzel~H is the sole conjectural input.  The conclusion remains
conditional: this argument neither proves Schinzel~H nor gives an
unconditional solution of Hilbert's Tenth Problem over $\Q$.
% src: data/l22_elimination.jsonl (fixed-Z-theorem row)
% src: data/l20_admissible.jsonl; data/l18_schinzel_implies_h.jsonl
\end{remark}
